\section{Durchführung und Ergebnisse}

\subsection{Vergleich der Teststrategien anhand spezifischer Anwendungsfälle}
Die beiden Teststrategien \texttt{Testpyramide} und \texttt{Honeycomb-Modell} wurden anhand praxisnaher Anwendungsfälle evaluiert. Die Analyse umfasst den Bewerbungsprozess, das Vertragsmanagement und das Eventmanagement. Die Messungen (Laufzeit, Anzahl durchgeführter Tests) sowie die Testabdeckung (vgl. Abschnitte \ref{tab:jacoco-ergebnisse} und \ref{tab:pitest-detailed-logs}) wurden in einer einheitlichen Testumgebung durchgeführt, um die Ergebnisse vergleichbar zu halten. 

\subsubsection{Bewerbungsprozess}
Im Bewerbungsprozess zeigte sich, dass die \texttt{Testpyramide} durch die Verwendung einer H2-Datenbank geringere Laufzeiten aufweist, während das \texttt{Honeycomb-Modell} mehr Zeit in Integrationspfade investiert und dabei MySQL-Testcontainer nutzt. Die Ergebnisse sind in Tabelle \ref{tab:bewerbungsprozess} dargestellt.

\begin{table}[h]
	\centering
	\caption{Laufzeiten und Testanzahl im Bewerbungsprozess}
	\label{tab:bewerbungsprozess}
	\begin{tabular}{|l|l|l|l|l|}
		\hline
		\textbf{Modell} & \textbf{Test Typ} & \textbf{Laufzeit} & \textbf{Anzahl der Tests} & \textbf{Datenbank} \\
		\hline
		Honeycomb       & Integration       & \textbf{1 s 582 ms} & 18  & Testcontainer MySQL \\
		\hline
		Testpyramide    & Alle              & \textbf{1 s 876 ms} & 44  & H2 \\
		\hline
		Testpyramide    & Alle              & \textbf{3 s 408 ms} & 44  & Testcontainer MySQL \\
		\hline
	\end{tabular}
\end{table}

\subsubsection{Vertragsmanagement}
Das Vertragsmanagement, als zentraler und vergleichsweise aufwendiger Anwendungsfall, zeigt deutliche Unterschiede zwischen den Strategien. Tabelle \ref{tab:vertragsmanagement} veranschaulicht, dass das \texttt{Honeycomb-Modell} infolge umfangreicher Integrationstests mit MySQL länger läuft, während die \texttt{Testpyramide} kürzere Zeiten bei einer größeren Testanzahl verzeichnet.

\begin{table}[h]
	\centering
	\caption{Laufzeiten und Testanzahl im Vertragsmanagement}
	\label{tab:vertragsmanagement}
	\begin{tabular}{|l|l|l|l|l|}
		\hline
		\textbf{Modell} & \textbf{Test Typ} & \textbf{Laufzeit} & \textbf{Anzahl der Tests} & \textbf{Datenbank} \\
		\hline
		Honeycomb       & Integration       & \textbf{9 s 835 ms} & 48  & Testcontainer MySQL \\
		\hline
		Testpyramide    & Alle              & \textbf{2 s 565 ms} & 103 & H2 \\
		\hline
		Testpyramide    & Alle              & \textbf{3 s 757 ms} & 103 & Testcontainer MySQL \\
		\hline
	\end{tabular}
\end{table}

\subsubsection{Eventmanagement}
Beim Eventmanagement werden komplexe Datenbankvorgänge und mehrere miteinander interagierende Module getestet. Tabelle \ref{tab:eventmanagement} belegt, dass das \texttt{Honeycomb-Modell} im Integrationsmodus nochmals mehr Laufzeit benötigt als die \texttt{Testpyramide}, die wiederum mit einer H2-Datenbank besonders schnelle Durchläufe erzielt.

\begin{table}[h]
	\centering
	\caption{Laufzeiten und Testanzahl im Eventmanagement}
	\label{tab:eventmanagement}
	\begin{tabular}{|l|l|l|l|l|}
		\hline
		\textbf{Modell} & \textbf{Test Typ} & \textbf{Laufzeit} & \textbf{Anzahl der Tests} & \textbf{Datenbank} \\
		\hline
		Honeycomb       & Integration       & \textbf{12 s 939 ms} & 58  & Testcontainer MySQL \\
		\hline
		Testpyramide    & Alle              & \textbf{2 s 858 ms}  & 140 & H2 \\
		\hline
		Testpyramide    & Alle              & \textbf{7 s 428 ms}  & 140 & Testcontainer MySQL \\
		\hline
	\end{tabular}
\end{table}

\subsection{Zusammenfassung der Testabdeckung mit JaCoCo}
Um die Teststrategien \texttt{Testpyramide} und \texttt{Honeycomb-Modell} hinsichtlich ihrer Abdeckung zu bewerten, wurde das Tool \texttt{JaCoCo} eingesetzt. Es erfasst verschiedene Metriken (z.\,B. \texttt{Instructions}, \texttt{Branches}, \texttt{Methods} und \texttt{Classes}) und ermöglicht eine Differenzierung, wie viele Teile des Codes (in Zeilen und Verzweigungen) von den Tests durchlaufen werden. Die Resultate sind in Tabelle \ref{tab:jacoco-ergebnisse} zu sehen.

\begin{table}[h]
	\centering
	\caption{Testabdeckung mit JaCoCo für Honeycomb-Modell und Testpyramide}
	\label{tab:jacoco-ergebnisse}
	\begin{tabular}{|l|l|l|}
		\hline
		\textbf{Metrik}       & \textbf{Honeycomb} & \textbf{Testpyramide} \\ \hline
		Instructions          & 91 \%              & 94 \%                 \\ \hline
		Branches              & 65 \%              & 80 \%                 \\ \hline
		Lines                 & 91 \% (Missed: 51) & 94 \% (Missed: 35)    \\ \hline
		Methods               & 87.6 \%            & 91.1 \%               \\ \hline
		Classes               & 98 \%              & 100 \%                \\ \hline
	\end{tabular}
\end{table}

Die Messungen zeigen, dass beide Ansätze eine insgesamt hohe Codeabdeckung erzielen. Während die \texttt{Testpyramide} in den meisten Metriken leicht vorn liegt (z.\,B. \texttt{Branches}, \texttt{Methods}), stellt das \texttt{Honeycomb-Modell} mit weniger Tests eine ebenfalls hohe Abdeckung sicher, vor allem in \texttt{Instructions} (91\,\%) und \texttt{Classes} (98\,\%).

\subsection{Mutation Testing mit PITest}
Zur Beurteilung, wie gut die Tests nicht nur \emph{Code ausführen}, sondern auch \emph{Fehler aufdecken} können, kam \texttt{PITest} zum Einsatz. Dieses Werkzeug generiert Mutationen (z.\,B. veränderte Return-Werte oder Negation von Bedingungen), um zu prüfen, ob sie von den Tests entdeckt und „getötet“ werden.

\subsubsection{Ergebnisse der Mutation Tests}
Tabelle \ref{tab:pitest-detailed-logs} zeigt exemplarisch, wie viele Mutationen (\texttt{Gen.}) pro Mutator erzeugt, getötet (\texttt{Killed}) und u.\,U. ohne Abdeckung (\texttt{NoCov}) geblieben sind. Überlebende Mutationen (\texttt{Surv.}) geben Aufschluss über mögliche Lücken in den existierenden Tests. 

\begin{table}[h!]
	\centering
	\caption{Detaillierte PITest-Log-Daten im Vergleich (Honeycomb vs.\ Testpyramide)}
	\label{tab:pitest-detailed-logs}
	\resizebox{\textwidth}{!}{%
		\begin{tabular}{|l||c|c|c|c|c||c|c|c|c|c|}
			\hline
			\multirow{2}{*}{\textbf{Mutator}} 
			& \multicolumn{5}{c||}{\textbf{Honeycomb-Modell}} 
			& \multicolumn{5}{c|}{\textbf{Testpyramide}} \\
			\cline{2-11}
			& \textbf{Gen.} & \textbf{Killed(\%)} & \textbf{Surv.} & \textbf{TOut} & \textbf{NoCov} 
			& \textbf{Gen.} & \textbf{Killed(\%)} & \textbf{Surv.} & \textbf{TOut} & \textbf{NoCov} \\
			\hline
			\textbf{ConditionalsBoundary} & 2 & 2 (100\,\%) & 0 & 2 & 0 
			& 2 & 1 (50\,\%)  & 1 & 0 & 0 \\
			\textbf{PrimitiveReturns}     & 1 & 1 (100\,\%) & 0 & 1 & 0
			& 1 & 1 (100\,\%) & 0 & 0 & 0 \\
			\textbf{VoidMethodCall}       & 9 & 8 (89\,\%)  & 0 & 6 & 1
			& 9 & 8 (89\,\%)  & 1 & 0 & 0 \\
			\textbf{BooleanTrueReturnVals}& 4 & 4 (100\,\%) & 0 & 4 & 0
			& 4 & 4 (100\,\%) & 0 & 0 & 0 \\
			\textbf{NullReturnVals}       & 78 & 68 (87\,\%) & 0 & 61 & 10
			& 78 & 72 (92\,\%) & 0 & 3 & 6 \\
			\textbf{EmptyObjectReturnVals}& 19 & 19 (100\,\%) & 0 & 19 & 0
			& 19 & 18 (95\,\%)  & 0 & 2 & 1 \\
			\textbf{BooleanFalseReturnVals} & 5 & 5 (100\,\%) & 0 & 5 & 0
			& 5 & 4 (80\,\%)  & 1 & 0 & 0 \\
			\textbf{NegateConditionals}     & 33 & 31 (94\,\%) & 0 & 31 & 2
			& 33 & 33 (100\,\%)& 0 & 1 & 0 \\
			\hline
		\end{tabular}
	}
\end{table}

\paragraph{Zusammenfassung der wichtigsten Kennzahlen}
\begin{itemize}
	\item \textbf{Honeycomb-Modell:}
	\begin{itemize}
		\item Generierte Mutationen: 151
		\item Getötete Mutationen: 138 (91\,\%)
		\item Mutationen ohne Abdeckung (NO\_COVERAGE): 13
		\item Line Coverage (mutierte Klassen): 92\,\% (495/537)
		\item Test Strength: 100\,\% (alle abgedeckten Mutationen wurden getötet)
		\item Timeouts: 129
		\item Gesamtlaufzeit: ca. 52 Minuten 21 Sekunden
	\end{itemize}
	\item \textbf{Testpyramide:}
	\begin{itemize}
		\item Generierte Mutationen: 151
		\item Getötete Mutationen: 141 (93\,\%)
		\item Mutationen ohne Abdeckung (NO\_COVERAGE): 7
		\item Line Coverage (mutierte Klassen): 94\,\% (507/537)
		\item Test Strength: 98\,\% (141/144 abgedeckte Mutationen wurden getötet)
		\item Timeouts: 6
		\item Gesamtlaufzeit: ca. 5 Minuten 25 Sekunden
	\end{itemize}
\end{itemize}

\begin{table}[h]
	\centering
	\caption{Kompakte Gegenüberstellung der Hauptmetriken (Honeycomb vs.\ Testpyramide)}
	\label{tab:pittest-ergebnisse}
	\begin{tabular}{|l|c|c|}
		\hline
		\textbf{Metrik} & \textbf{Honeycomb} & \textbf{Testpyramide} \\
		\hline
		Line Coverage (mutierte Klassen) & 92\,\% (495/537) & 94\,\% (507/537) \\
		\hline
		Mutation Coverage & 91\,\% (138/151) & 93\,\% (141/151) \\
		\hline
		Mutationen ohne Abdeckung (NO\_COVERAGE) & 13 & 7 \\
		\hline
		Überlebende Mutationen (SURVIVED) & 0 & 3 \\
		\hline
		Test Strength & 100\,\% (138/138) & 98\,\% (141/144) \\
		\hline
		Timeouts & 129 & 6 \\
		\hline
		Tests pro Mutation & 0{,}01 & 1{,}08 \\
		\hline
		Gesamtlaufzeit & 52 Min 21 Sek & 5 Min 25 Sek \\
		\hline
	\end{tabular}
\end{table}

\subsubsection{Zusammenfassung der Mutation-Testergebnisse}
Die in Tabelle \ref{tab:pittest-ergebnisse} aufgeführten Metriken machen deutlich, dass beide Strategien eine hohe Mutation Coverage aufweisen. Das \texttt{Honeycomb-Modell} benötigt jedoch mehr Zeit (rund 52 Minuten) und verzeichnet 129 Timeouts. Die \texttt{Testpyramide} führt zwar mehr Einzeltests aus, schliesst ihre Durchläufe jedoch in knapp 5\,Minuten ab und weist deutlich weniger Timeouts auf. 

\subsection{Vergleich der Datenbanklösungen}
Der Einsatz einer H2-Datenbank führt zu \emph{kürzeren} Testlaufzeiten und ist in vielen Fällen für Unit- und (kleinere) Integrationstests ausreichend schnell. Testcontainer-basierte MySQL-Instanzen bilden hingegen \emph{produktionsnähere Bedingungen} ab, was jedoch in komplexeren Szenarien höhere Laufzeiten und potenziell mehr Timeouts nach sich zieht. Beide Ansätze können je nach Zielsetzung miteinander kombiniert werden, um ein ausgewogenes Verhältnis aus Geschwindigkeit und Realitätsnähe zu erreichen.

\subsection{Zusammenfassung der Ergebnisse}
Die durchgeführten Tests in den drei Anwendungsfällen – Bewerbungsprozess, Vertragsmanagement und Eventmanagement – zeigen, dass das \texttt{Honeycomb-Modell} und die \texttt{Testpyramide} jeweils charakteristische Laufzeiten, Testumfänge und Abdeckungswerte aufweisen. Anhand der JaCoCo-Metriken (Tabelle \ref{tab:jacoco-ergebnisse}) ist erkennbar, dass beide Strategien hohe Codeabdeckung liefern, wobei die \texttt{Testpyramide} in den meisten Kennzahlen (z.\,B. Branch-Coverage) noch etwas weiter reicht. Die Messungen mit \texttt{PITest} (Tabellen \ref{tab:pitest-detailed-logs} und \ref{tab:pittest-ergebnisse}) bestätigen, dass beide Ansätze Mutationen effektiv erkennen und töten können, wenngleich sich \emph{Laufzeit} und \emph{Anzahl der Timeouts} deutlich unterscheiden. 

H2 ermöglicht bei vielen Tests eine sehr kurze Durchlaufzeit, während MySQL-Testcontainer realistischere Datenbankszenarien und damit umfassendere Integrationsprüfungen erlauben. In allen drei Fällen (Bewerbungsprozess, Vertragsmanagement, Eventmanagement) waren die Laufzeiten beim \texttt{Honeycomb-Modell} länger, häufig bei vergleichsweise wenigen Tests. Die \texttt{Testpyramide} zeigte mehr Einzeltests und kürzere Gesamtdauer, insbesondere mit H2. Somit bleibt festzuhalten, dass beide Strategien je nach Einsatzzweck – \emph{Realitätsnähe} versus \emph{Schnelligkeit und Detailtiefe} – sinnvoll eingesetzt werden können.
